\subsection{重数与有限扩张}

命题 5. — 假设 A 是半局部环，且 \(\rho: A \to B\) 是使 B 成为有限生成 A-模的环同态。设 N 是一个非零有限生成的 B-模，q 是 A 的一个理想，包含于 A 的根基中，并且 N/qN 具有有限长度。在 B 的极大理想中，将满足条件的那些记为 \(\mathfrak{m}_1, \ldots, \mathfrak{m}_r\)，其中 \(\dim_{\mathrm{B}_{\mathfrak{m}_i}} (\mathrm{N}_{\mathfrak{m}_i}) = \dim_{\mathrm{B}} (\mathrm{N})\)（根据 IV, § 2, 第 5 号，命题 9 的推论 3，它们的数目有限）。置 \(\mathbf{B}_i = \mathbf{B}_{\mathfrak{m}_i}\)，\(\mathfrak{q}_i = \mathfrak{q}\mathbf{B}_i\)，其中 \(1 \leqslant i \leqslant r\)。则有以下等式

\[
\dim_ {A} (N) = \dim_ {B} (N),
\]

\[
e _ {\mathfrak {q} \mathrm{B}} ^ {\mathbf {B}} (\mathbf {N}) = \sum_ {i = 1} ^ {r} e _ {\mathfrak {q} _ {i}} ^ {\mathbf {B} _ {i}} (\mathbf {N} _ {\mathfrak {m} _ {i}}),
\]

\[
e _ {\mathfrak {q}} ^ {\mathrm{A}} (\mathrm{N}) = \sum_ {i = 1} ^ {r} \left[ \mathrm{B} / \mathfrak {m} _ {i}: \mathrm{A} / \rho^ {- 1} (\mathfrak {m} _ {i}) \right] e _ {\mathfrak {q} _ {i}} ^ {\mathrm{B} _ {i}} (\mathrm{N} _ {\mathfrak {m} _ {i}}).
\]

第一个等式由 § 2，第 3 号，定理 1 c) 得出；第二个等式由第 1 号的注记 4 得出（注意，m \(_{i}\) 属于 V(qB)（\(1 \leqslant i \leqslant r\)），因为根据 V, § 2，第 1 号，命题 1，有 \(\rho^{-1}(m_{i}) \supset q\)）。下面证明第三个等式。设 E 是一个有限长度的 B-模；有

\[
\operatorname{long} _ {\mathbf {A}} (\mathbf {E}) = \sum_ {m} \left[ \mathrm{B} / m: \mathrm{A} / \rho^ {- 1} (m) \right]. \operatorname{long} _ {\mathrm{B} _ {m}} \left(\mathrm{E} _ {m}\right),
\]

m 遍历 B 的极大理想集合；当 E 是某个 B/m 时，这一点确实显然；一般情形由此得到，因为 E 有一个合成序列，其商同构于 B/m。将这个公式应用于 B-模 N/q \(^{n+1}\) N，便根据重数的定义得到所需等式。

推论。— 若对于所有 i 都有 \(\left[\mathrm{B}/\mathfrak{m}_{i}:\mathrm{A}/\rho^{-1}(\mathfrak{m}_{i})\right]=1\)，则有 \(e_{\mathfrak{q}}^{\mathrm{A}}(\mathrm{N})=e_{\mathfrak{q}\mathrm{B}}^{\mathrm{B}}(\mathrm{N})\)。

引理 1. — 设 \(\rho: A \to B\) 是环同态，且 \(\mathfrak{p}\) 是 A 的一个素理想。考虑以下两个性质：

\begin{enumerate}
\def\labelenumi{(\roman{enumi})}
\item
  典范同态 \(\tilde{\rho}\) 从 \(\mathbf{A}_{\mathfrak{p}}\) 到 \(\mathbf{A}_{\mathfrak{p}} \otimes_{\mathbf{A}} \mathbf{B}\) 是双射；
\item
  存在 B 中唯一的素理想 \(\mathfrak{r}\)，它位于 \(\mathfrak{p}\) 之上，且典范同态 \(\rho_{\mathfrak{p}}\) 从 \(A_{\mathfrak{p}}\) 到 \(B_{\mathfrak{r}}\) 是双射。
\end{enumerate}

我们有 \((\mathrm{i}) \Rightarrow (\mathrm{ii})\)。若 \(\mathfrak{p}\) 是极小的，或者 \(\mathbf{B}\) 整于 \(\mathbf{A}\)，则有 \((\mathrm{i}) \Leftrightarrow (\mathrm{ii})^{1}\)。

环 \(\mathbf{A}_{\mathfrak{p}} \otimes_{\mathbf{A}} \mathbf{B}\) 可识别为 B 的分式环 \(S^{-1} \mathbf{B}\)，其中乘法子集 \(S = \rho(A - \mathfrak{p})\) 是 B 的子集。因而 \(S^{-1} \mathbf{B}\) 的素理想正是 \(S^{-1} \mathfrak{q}\)，其中 \(\mathfrak{q}\) 是满足 \(\rho^{-1}(\mathfrak{q}) \subset \mathfrak{p}\) 的 B 的素理想；若 \(\mathfrak{q}\) 是这样的理想，则 \((S^{-1} \mathbf{B})_{S^{-1} \mathfrak{q}}\) 可识别为 \(\mathbf{B}_{\mathfrak{q}}\)（II, § 2, 第 5 号，命题 11）。

若性质 (i) 成立，则存在（V, § 2，第 1 号，引理 1）中所述的 B 中唯一的素理想 \(\mathfrak{r}\)，满足 \(\rho^{-1}(\mathfrak{r}) = \mathfrak{p}\)。此外，\(B_{\mathfrak{r}}\) 可识别为分式环 \((S^{-1}B)_{S^{-1}\mathfrak{r}}\)，因而也可识别为 \((A_{\mathfrak{p}})_{s}\)，其中 \(s\) 是 \(S^{-1}\mathfrak{r}\) 在同构 \(\tilde{\rho}: A_{\mathfrak{p}} \to S^{-1}B\) 下的逆像；并且 \(\tilde{\rho}^{-1}(S^{-1}\mathfrak{r}) = (A - \mathfrak{p})^{-1}\mathfrak{p} = \mathfrak{p}A_{\mathfrak{p}}\)，从而有 (ii)。

反之，若 (ii) 成立，令 r 为位于 p 之上的 B 的唯一素理想。由于 \((\mathrm{S}^{-1}\mathrm{B})_{\mathrm{S}^{-1}\mathrm{r}}\) 可识别为 \(B_{r}\)，只需证明 \(S^{-1}B\) 是以 \(S^{-1}r\) 为极大理想的局部环；也就是说，B 的每个满足 \(\rho^{-1}(q) \subset p\) 的素理想 q 都包含于 r。若 p 是极小的，则有 \(\rho^{-1}(q) = p\)，故 q = r。若 B 整于 A，则根据 V, § 2，第 1 号，定理 1 的推论 2，存在 B 的一个素理想 \(r'\)，使 \(q \subset r'\) 且 \(\rho^{-1}(r') = p\)；必有 \(r' = r\)，故 \(q \subset r\)。

引理 2.--- 假设 A 是半局部环；设 q 是 A 的一个定义理想，且 \(\rho : A \to B\) 是使 B 成为有限生成 A-模的环同态。假设，对于 A 的每个素理想 \(\mathfrak{p}\)（必为极小的），若满足 \(\dim(A/\mathfrak{p}) = \dim(A)\)，则存在 B 中唯一的素理想 \(\mathfrak{r}\)，位于 \(\mathfrak{p}\) 之上，且典范同态 \(\rho_{\mathfrak{p}} : A_{\mathfrak{p}} \to B_{\mathfrak{r}}\) 是双射。则有 \(\dim_{A}(B) = \dim(A)\) 且 \(e_{\mathfrak{q}}^{A}(B) = e_{\mathfrak{q}}^{A}(A)\)。记 \(\mathfrak{S}_{A}\)（分别为 \(\mathfrak{S}_{B}\)）为如下集合：A 的素理想 \(\mathfrak{p}\) 满足

\[
\dim (\mathrm{A} / \mathfrak {p}) = \dim (\mathrm{A}) (\text { resp. } \dim_ {\mathrm{A}} (\mathrm{B} / \mathfrak {p} \mathrm{B}) = \dim_ {\mathrm{A}} (\mathrm{B}));
\]

我们有 \(\mathfrak{S}_{\mathbf{A}} \neq \emptyset\)。取 \(\mathfrak{p} \in \mathfrak{S}_{\mathbf{A}}\)；由假设，存在一个位于 \(\mathfrak{p}\) 之上的 B 的素理想。于是 \(\rho^{-1}(\mathfrak{pB}) = \mathfrak{p}\)（II, § 2, 第 5 号，命题 11 的推论 3），并且

\[
\dim (\mathrm{A} / \mathfrak {p}) = \dim (\mathrm{B} / \mathfrak {p B}) = \dim_ {\mathrm{A}} (\mathrm{B} / \mathfrak {p B})
\]

根据 § 2，第 3 号，定理 1 b) 和 c)，因此有

\[
\dim_ {\mathrm{A}} (\mathrm{B}) \geqslant \dim_ {\mathrm{A}} (\mathrm{B} / \mathfrak {p B}) = \dim (\mathrm{A} / \mathfrak {p}) = \dim (\mathrm{A}) \geqslant \dim_ {\mathrm{A}} (\mathrm{B}).
\]

这说明 \(\mathfrak{S}_{\mathrm{A}} \subset \mathfrak{S}_{\mathrm{B}}\) 且 \(\dim(\mathrm{A}) = \dim_{\mathrm{A}}(\mathrm{B})\)。反之，若 \(\mathfrak{p} \in \mathfrak{S}_{\mathrm{B}}\)，则有不等式

\[
\dim_ {A} (B / p B) = \dim_ {A} (B) = \dim (A) \geqslant \dim (A / p) \geqslant \dim_ {A} (B / p B),
\]

从而 \(p \in \mathfrak{S}_A\)，且 \(\mathfrak{S}_B = \mathfrak{S}_A\)。根据第 1 号的命题 3 及其推论，有

\[
e _ {\mathfrak {q}} ^ {\mathrm{A}} (\mathrm{A}) = \sum_ {\mathfrak {p} \in \mathfrak {S} _ {\mathrm{A}}} e _ {\mathfrak {q}} ^ {\mathrm{A}} (\mathrm{A} / \mathfrak {p}) \quad \text { and } \quad e _ {\mathfrak {q}} ^ {\mathrm{A}} (\mathrm{B}) = \sum_ {\mathfrak {p} \in \mathfrak {S} _ {\mathrm{B}}} \operatorname{long} _ {\mathrm{A} _ {\mathfrak {p}}} (\mathrm{A} _ {\mathfrak {p}} \otimes_ {\mathrm{A}} \mathrm{B}) e _ {\mathfrak {q}} ^ {\mathrm{A}} (\mathrm{A} / \mathfrak {p});
\]

根据引理 1，有 \(\operatorname{long}_{A_{\mathfrak{p}}}(A_{\mathfrak{p}} \otimes_{A} B) = 1\)，其中 \(\mathfrak{p} \in \mathfrak{S}_{A}\)；从而 \(e_{\mathfrak{q}}^{A}(A) = e_{\mathfrak{q}}^{A}(B)\)。

命题 6. --- 假设 A 是半局部且约化的；设 q 是 A 的一个定义理想；设 A' 是 A 的全分式环，B 是 A' 的一个有限 A-子代数。则 B 是半局部的，且 qB 是它的一个定义理想。假设，对于 B 的每个满足 \(\dim(\mathbf{B}_{\mathfrak{m}})=\dim(\mathbf{B})\) 的极大理想 m，都有 \([\mathbf{B}/\mathfrak{m}:\mathbf{A}/(\mathbf{A}\cap\mathfrak{m})]=1\)。则有 \(e_{\mathfrak{q}}^{\mathbf{A}}(\mathbf{A})=e_{\mathfrak{q}\mathbf{B}}^{\mathbf{B}}(\mathbf{B})\)。

根据 IV, § 2，第 5 号，命题 9 的推论 3，B 是半局部的，定义理想为 qB。根据命题 5 的推论，有 \(e_{qB}^{B}(B) = e_{q}^{A}(B)\)。由于 A' 可识别为 \(\prod_{p} A_p\)，其中 p 遍历 A 的极小素理想集合（IV, § 2，第 5 号，命题 10），故对于 A 的每个极小素理想 p，典范映射 \(A_{p} \to A_{p} \otimes_{A} B\) 都是双射。于是由引理 1 和引理 2，\(e_{q}^{A}(B) = e_{q}^{A}(A)\)，从而命题得证。

例。--- 设 \(k\) 是特征 \(\neq 2\) 的域，取 A 为局部环 \(k[[X, Y]]/(X^{2} + Y^{2})\)，其剩余域为 \(k\)。取 \(B = k[[X, T]]/(T^{2} + 1)\)，其中 \(T = Y/X\)。分两种情形：若 -1 是元素 \(i\)（来自 \(k\)）的平方，则 B 有两个极大理想，分别由 \(\{X, T + i\}\) 和 \(\{X, T - i\}\) 生成；它们的剩余域为 \(k\)，并且有 \(e_{\mathfrak{m}_{A}}^{A}(A) = e_{\mathfrak{m}_{A}B}^{B}(B) = 2\)。若 -1 不是 \(k\) 中的平方，则 B 有唯一极大理想 (X)，其剩余域为 \(k[T]/(T^{2} + 1)\)，且有 \(e_{\mathfrak{m}_{A}}^{A}(A) = 2\)，\(e_{\mathfrak{m}_{A}B}^{B}(B) = 1\)。

